\chapter{La règle générale}

Le chapitre précédent comptait des cailloux. Celui-ci écrit la même chose en
symboles, et ajoute la seule subtilité de toute la méthode.

\section{Ce qu'on calcule vraiment : une intégrale}

Une surface, une quantité de lumière, une moyenne sur toutes les parties
possibles : ce sont toutes des \textbf{intégrales}. Une intégrale, c'est une
somme d'une infinité de valeurs sur un domaine :

\[ I = \int_{D} f(x)\,dx \]

Lisez-la ainsi : « pour chaque point $x$ du domaine $D$, prends la valeur
$f(x)$, et fais-en la somme ». Le problème est qu'il y a une infinité de
points — on ne peut pas les visiter.

\section{L'estimateur}

Monte-Carlo remplace cette somme infinie par une moyenne finie :

\begin{nkretenir}[La formule centrale, à retenir par cœur]
\[ \langle I \rangle_N \;=\; \frac{1}{N}\sum_{i=1}^{N} \frac{f(x_i)}{p(x_i)} \]

où les $x_i$ sont $N$ points tirés au hasard selon une densité de probabilité
$p$. Ce nombre s'approche de $I$ quand $N$ grandit.
\end{nkretenir}

Le $f(x_i)$ ne pose aucune difficulté : c'est la valeur de ce qu'on somme, au
point tiré. Le $\frac{1}{N}\sum$ non plus : c'est une moyenne. Reste le
$p(x_i)$ au dénominateur, et c'est lui que tout le monde rate.

\section{Pourquoi diviser par $p$ — le point clé}

Imaginez que vous comptiez les arbres d'une forêt en vous y promenant, puis que
vous multipliiez par la surface. Si votre promenade passe deux fois plus souvent
près de la rivière que sur la colline, vous compterez les arbres de la rivière
deux fois trop. Votre estimation sera fausse — non parce que vous avez mal
compté, mais parce que \textbf{votre façon de vous promener favorisait un
endroit}.

La correction est évidente une fois dite : chaque arbre vu doit être divisé par
la fréquence avec laquelle vous passiez là. Un arbre vu dans une zone
sur-visitée compte moins ; un arbre vu dans une zone rarement visitée compte
davantage.

\begin{nkretenir}[Le rôle de $p$]
La division par $p$ \textbf{annule le favoritisme de votre tirage}. Elle
transforme n'importe quelle façon de tirer — même très biaisée — en une
estimation correcte.

C'est ce qui rend la méthode libre : vous pouvez tirer \emph{où vous voulez},
du moment que vous savez avec quelle probabilité vous y tirez.
\end{nkretenir}

\subsection{Le cas uniforme, que vous connaissez déjà}

Si vous tirez uniformément sur un domaine de volume $V$, alors la densité vaut
$p(x) = 1/V$ partout — la même chance partout. Remplacez dans la formule :

\[ \langle I \rangle_N = \frac{1}{N}\sum_{i=1}^{N} \frac{f(x_i)}{1/V}
   = \frac{V}{N}\sum_{i=1}^{N} f(x_i) \]

C'est-à-dire : \textbf{le volume, multiplié par la moyenne des valeurs}. C'est
exactement le calcul de la flaque (surface de la cour $\times$ proportion de
plouf) et celui de $\pi$ (surface du carré $\times$ proportion dans le cercle).
Vous utilisiez déjà la formule générale sans le savoir.

\section{Non biaisé : viser juste sans être précis}

Une propriété mérite un mot, parce qu'elle est contre-intuitive.

L'estimateur est dit \textbf{non biaisé} : son espérance est exactement $I$,
\emph{dès le premier échantillon}. Avec un seul caillou, votre estimation de la
flaque sera ridicule (0 \% ou 100 \%) — mais elle n'est pas \emph{décalée} :
si vous refaisiez l'expérience à un caillou des millions de fois, la moyenne de
toutes ces estimations ridicules donnerait la bonne surface.

\begin{nkretenir}[Justesse et précision ne sont pas la même chose]
\begin{itemize}
  \item \textbf{Non biaisé} = qui vise le centre de la cible. Les impacts sont
        dispersés, mais leur centre est le bon.
  \item \textbf{Précis} = qui groupe ses impacts. Ils peuvent être groupés à
        côté de la cible.
\end{itemize}
Monte-Carlo est \emph{juste} d'emblée et devient \emph{précis} avec le temps.
Un estimateur biaisé, lui, est stable et faux : il ne bruite pas, il ment. Et
c'est bien plus grave, parce que rien ne le signale.
\end{nkretenir}

\begin{nkpiege}[Deux conditions à ne jamais oublier]
La formule n'est valable que si :
\begin{enumerate}
  \item $p(x) > 0$ \textbf{partout où $f(x) \neq 0$}. Si vous ne tirez jamais
        dans une région où la fonction vaut quelque chose, cette contribution
        est perdue pour toujours — et aucun nombre d'échantillons ne la
        rattrapera. En rendu, c'est la lumière qu'on ne verra jamais parce
        qu'aucun rayon ne va dans sa direction ;
  \item $p$ est une \textbf{vraie} densité : positive, et d'intégrale 1. Une
        densité mal normalisée décale silencieusement tout le résultat.
\end{enumerate}
\end{nkpiege}

\begin{nkexo}[Vérifier que la division par $p$ est indispensable]
Estimez $\int_0^1 x^2\,dx$ (la réponse exacte est $1/3$) de deux façons :

\begin{enumerate}
  \item en tirant $x$ uniformément dans $[0,1]$, donc $p(x) = 1$ ;
  \item en tirant $x = \sqrt{u}$ avec $u$ uniforme — ce qui favorise les
        grandes valeurs de $x$. La densité correspondante est $p(x) = 2x$.
\end{enumerate}

Faites le second calcul \textbf{en oubliant volontairement} de diviser par
$2x$ : vous trouverez $1/2$ au lieu de $1/3$. Puis divisez : vous retrouvez
$1/3$. L'erreur de 50 \% ne vient pas du hasard — elle vient du favoritisme non
corrigé.
\end{nkexo}
